\documentclass[11pt]{article}
\usepackage[a4paper,margin=22mm]{geometry}
\usepackage{amsmath,amssymb,amsthm,mathtools,bm}
\usepackage{booktabs,array,microtype,xurl,hyperref}
\hypersetup{colorlinks=true,linkcolor=blue,citecolor=blue,urlcolor=blue,
pdftitle={P54 covariant EFT, finite Yukawa operators, and sequential seesaw}}
\setlength{\emergencystretch}{2em}
\setcounter{tocdepth}{1}
\newcommand{\tr}{\operatorname{Tr}}
\newcommand{\MS}{\overline{\mathrm{MS}}}
\newcommand{\diag}{\operatorname{diag}}
\newcommand{\dd}{\mathrm d}
\newcommand{\B}{\mathcal B}
\newtheorem{proposition}{Proposition}
\title{Route-F P54: Covariant Lower Action,\\Finite Yukawa Operators,\\and Moving Seesaw Thresholds}
\author{P54PQ-v2 four-dimensional mainline}
\date{6 September 2026}
\begin{document}
\maketitle
\begin{abstract}
The previous common-phase/PS-flow checkpoint was committed and pushed as
\texttt{4ad74e7}. Keeping that same action fixed, we construct a local
PS-covariant scalar/vector effective action and compute actual finite
pure-Yukawa diagrams at the upper saddle. Split real six-plet masses
generate conjugate-bidoublet operators outside the former four-matrix
boundary. Their coefficients are calculable functions of the original two
family matrices, not additional fitting parameters. Independently, we
compute the full four-doublet scalar-cubic kinetic matrix and implement
moving, possibly crossing or degenerate, sterile-neutrino thresholds.
The latter includes the corrected dimension-five feedback found in 2024.
These are bounded matching results, not a completed gauge/ghost finite
matching calculation, physical flavor fit, or Higgs pole prediction.
\end{abstract}
\tableofcontents
\clearpage

\section{Frozen action and approximation ledger}
The field content remains $3\,16_F+54_{H,\mathbb R}+126_{H,\mathbb C}
+10_{H,\mathbb C}+1_{H,\mathbb C}$ with the action and common phase
dictionary in the preceding P54 cards. There are still only two symmetric
UV family spurions $h=h_{\rm raw}$ and $f=f_{\rm raw}$. All scalar masses
and vertices below come from that same 328-real-coordinate polynomial.
No new vacuum or lattice scan, field multiplet, free mirror matrix, or
physical scale is fitted here.

The actual broken background has $(\omega,\sigma,v_s)=(1,0.1265,0.25)$.
The nearby PS saddle has $(1.00081030592151,0,0.254277697083708)$.
Historical dimensionful unification scales remain unrecertified. A
synthetic family benchmark demonstrates covariance and implementation,
not agreement with measured masses or mixings.

There are three distinct approximations: (i) a tree local two-derivative
PS action after upper-field elimination; (ii) finite one-loop diagram
subsets with frozen tree propagators; (iii) one-loop dimension-five
running with tree threshold matching. They cannot simply be summed into
a claimed two-site one-loop prediction. The appropriate common
hard-minus-EFT subtraction is still required.

\section{A covariant lower action, not parent-weighted logarithms}
\subsection{Local scalar valley}
Let $B_A$ be the raw 248-real-dimensional PS parent embedding and $B_H$
the 55-dimensional upper physical scalar embedding. In a local slice,
\begin{equation}
 X(y)=x_0+B_Ay+B_Hz(y),\qquad B_H^T\nabla V(X(y))=0.
\end{equation}
At $y=0$, define $H=V^{(2)}(x_0)$ and
\begin{align}
 C&=B_H^THB_H,& W&=B_H^THB_A,& D&=-C^{-1}W,\\
 T&=X_{,y}=B_A+B_HD,&
 X_{,ab}&=-B_HC^{-1}B_H^TV^{(3)}[T_a,T_b].
\end{align}
The measured $\lambda_{\min}(C)=0.099450257411\,\omega^2>0$ proves
existence of an analytic valley germ by the implicit-function theorem.
It proves neither global continuation nor uniform scale separation.
Its local potential Hessian is $S=T^THT$.

\subsection{The 24-vector response}
For the real generators $R_A=-iT_A$, the convention $D=\partial-igAT$
means $D=\partial+gAR$. Write $O_A=R_AX$, $I=O^TO$ and
$j=gO^TD_{\rm PS}X$. The scalar kinetic term is
\begin{equation}
 \frac12\lVert D_{\rm PS}X+gOA_U\rVert^2.
\end{equation}
Its algebraic heavy-vector solution and substitution are
\begin{align}
 A_U^*&=-g^{-1}I^{-1}O^TD_{\rm PS}X,\\
 \Delta\mathcal L_2&=-\frac12j^T(g^2I)^{-1}j,\\
 \mathcal L_{\rm EFT}^{(2)}&=\frac12(Dy)^Tg(y)(Dy)-V(X(y)),\qquad
 g=T^T(1-OI^{-1}O^T)T. \label{eq:metric}
\end{align}
The 24 directions are a symmetric coset, not a gauge subgroup:
$[\mathfrak p,\mathfrak u]\subset\mathfrak u$ and
$[\mathfrak u,\mathfrak u]\subset\mathfrak p\ne0$. Thus
$I^{-1}O^TT$ is a PS-covariant response one-form, not a principal
connection for an invented 24-dimensional group. The axion is held fixed
as a gauge-neutral spectator in this chart.

For a direction $q$, with $O_q=(R_AX_q)_A$,
\begin{align}
 I_q&=O_q^TO+O^TO_q,& K&=O^TT,& K_q&=O_q^TT+O^TT_q,\\
 a&=I^{-1}K,& a_q&=I^{-1}(K_q-I_qa),\\
 g_q&=T_q^TT+T^TT_q-K_q^Ta-K^Ta_q.
\end{align}
These supply actual derivative vertices, not merely a mass census.
The previous horizontal Schur metric is reproduced to $1.1\times10^{-15}$;
its eigenvalue interval is $[0.984249788817,1.01800727050]$.
All nine lower gauge masses and Goldstone normalizations agree with
the full action. The current-current subtraction has rank 12 and trace
$0.189002534197$.

\subsection{Nonlinear vertices and limits of the derivative expansion}
For the exported corrected neutral ray $q$, put
$J_H=B_H^TV^{(3)}[q,q]$. With $H^0=h/\sqrt2$,
\begin{equation}
 \lambda_{\rm upper}=\frac16V^{(4)}[q^4]
 -\frac12J_H^TC^{-1}J_H=0.407929900088,
 \quad \lambda_{\rm direct}=0.427823639280.
\end{equation}
The induced derivative term is
$0.024787549073\,h^2(\partial h)^2/\omega^2$.
These are the projected tree potential and metric on a bosonic-improved
ray, not an SM quartic or a pole mass. The retained, canonically
normalized 126 triplet has $m^2=0.268832826146\,\omega^2$ and
$hh$ source $0.00252743109350\,\omega$; eliminating it reconstructs
the full mixed 54/126 tree response.

Already at quadratic order a heavy scalar gives the exact Euclidean
kernel, in a raw scalar chart,
\begin{equation}
 K_E(p^2)=p^2G_0+A-W^T(C+p^2I)^{-1}W
 =S+p^2G-p^4W^TC^{-3}W+\cdots. \label{eq:resolvent}
\end{equation}
The corresponding exact pole calculation uses analytic continuation,
not a replacement of all generalized eigenvalues by exact masses.
The Yang--Mills term likewise yields higher-derivative operators upon
substituting $A_U^*$ into $F_U$ and $F_{\rm PS}$.
Only modes coupled through $W$ need this resolvent test; comparing every
retained mass to a single global minimum heavy mass would impose an
unnecessary restriction. The public API retains the matrices needed
to handle the genuinely coupled modes exactly.
For the twelve near-crossing modes, a source-selected invariant block is
\begin{equation}
 H_2/\omega^2=\begin{pmatrix}.2688378&.0013417351173\\
 .0013417351173&.261600225\end{pmatrix}.
\end{equation}
Its exact eigenvalues are $.261359494901$ and $.269078530099$.
This is a finite resolvent repair, not a reason to redesign the action.

\section{Actual finite upper Yukawa diagrams}
\subsection{Kernels, signs and scale consistency}
Use $\mathcal L_Y=-Y^a_{IJ}\psi_I\psi_Js_a/2+\mathrm{h.c.}$
with real canonical scalars and symmetric Weyl matrices. At the PS
saddle all 48 Weyl fermions are massless. For a positive heavy scalar
mass $M_b^2$ define
\begin{align}
 f_b&=-\frac1{16\pi^2}\int_0^1\log\frac{(1-t)M_b^2}{\mu^2}\dd t
 =\frac{1-\log(M_b^2/\mu^2)}{16\pi^2},\\
 h_b&=\frac1{16\pi^2}\int_0^1(1-t)
 \log\frac{(1-t)M_b^2}{\mu^2}\dd t
 =\frac{\frac12\log(M_b^2/\mu^2)-\frac14}{16\pi^2}.
\end{align}
The fermion self-energy shift supplies the factor $(1-t)$; the two
massless Weyl numerators in the triangle cancel one denominator.
Differentiating the two fermions cancels the action's symmetry factor.
The finite result, with $Z_\psi=1+K_\psi$, is
\begin{align}
 K_\psi&=-\sum_{b\in H}Y_b^\dagger Y_bh_b,\\
 \delta Y_a^{\rm 1PI}&=-\sum_{b\in H}Y_bY_a^\dagger Y_bf_b,\\
 \Delta Y_a&=\delta Y_a^{\rm 1PI}
 -\tfrac12(K_\psi^TY_a+Y_aK_\psi). \label{eq:finiteY}
\end{align}
The tensor normalization is computed from the actual Clifford
intertwiners, not imported from a different GUT. As an independent test,
\begin{equation}
 \frac{\dd\Delta Y_a}{\dd\log\mu}=-\frac1{16\pi^2}
 \left[\frac12(S_H^TY_a+Y_aS_H)
 +2\sum_bY_bY_a^\dagger Y_b\right],\quad S_H=\sum_bY_b^\dagger Y_b.
\end{equation}
This equals minus the actual UV-minus-PS beta contraction. All-massless
fermion contributions to scalar wave functions cancel in the upper
hard subtraction. A scalar-cubic/two-Yukawa triangle has no dimension-four
zero-momentum contribution with a massless internal fermion; it requires
a chirality flip or external momentum. This does not eliminate scalar
cubic wave-function diagrams.

\subsection{Finite conjugate-bidoublet operators}
Only 24 of the 55 heavy real scalars have nonzero Yukawa couplings:
\begin{center}
\begin{tabular}{lrr}
\toprule
Parent & Squared masses / $\omega^2$ & Multiplicities\\
\midrule
$10$ six & $0.190903147974,\ 0.200245831815$ & $6,6$\\
$126$ six & $0.372275090654,\ 0.415462413403$ & $6,6$\\
\bottomrule
\end{tabular}
\end{center}
For a complete degenerate pair $Y_I=\pm iY_R$,
$Y_RY_a^\dagger Y_R+Y_IY_a^\dagger Y_I=0$.
Unequal real masses weight these terms differently. Hence this
cancellation fails by an exact split logarithm, not roundoff.

Write $\Delta_x=\log(m_{x,+}^2/m_{x,-}^2)/(16\pi^2)$.
In the exported normalized mirror-intertwiner phase convention,
\begin{align}
 \widetilde H&=i\bigl(3\sqrt2\Delta_h\,hh^\dagger h
 -6\sqrt2\Delta_f\,fh^\dagger f\bigr),\\
 \widetilde F&=i\bigl(-4\Delta_h\,hf^\dagger h
 +8\Delta_f\,ff^\dagger f\bigr). \label{eq:mirror}
\end{align}
Their four real coefficients, in the displayed order, are
\begin{equation}
 (0.00128368533616,\ -0.00589775824991,\
 -0.00121027014148,\ 0.00556045980308).
\end{equation}
The fermion metric is $K_\psi=I_{16}\otimes k$, where
\begin{equation}
 k=0.0186815380325\,h^\dagger h+0.0108097087894\,f^\dagger f
\end{equation}
at $\mu/\omega=0.563058770439$.

The holomorphic projection still satisfies $F=\sqrt2R$ and $L=R$ in
this subset. But it is not the whole finite boundary. The mirror
matrices are orthogonal intertwiner directions and must be retained.
They are nonlinear functions of the same $h,f$, not new UV parameters.
The displayed $i$ is a tensor-basis phase. Restoring the PQ orbit gives
the required $e^{-4i\theta}$ spurion phase; there is no explicit PQ
breaking. Mirror terms vanish in the degenerate limit and are matching
scale independent.

A strict one-loop matching-plus-running expansion need not feed an
already one-loop mirror threshold back into a one-loop beta function:
that effect starts at two loops. A resummed flow retaining mirror
matrices must, however, evolve the complete real-scalar Weyl tensor.
It cannot use the former four-holomorphic-parent equations as a closed
system. In particular the paired-holomorphic fast cancellation cannot
be assumed for a scalar with both orientations of Yukawa coupling.

The six invariant families $H,F,L,R,\widetilde H,\widetilde F$ do close
under the generic one-loop real-Weyl contraction. The four bidoublet
matrices are general $3\times3$ matrices; $L,R$ are symmetric, giving
$4\cdot9+2\cdot6=48$ complex tensor components, not 48 new UV parameters.
PS representation decomposition excludes further dimension-four
Yukawa tensors. Both generic six-family and actual finite-boundary
tests reconstruct the complete beta tensor with relative error below
$5\times10^{-15}$. Restriction to four families leaves a nonzero
$3.09\times10^{-4}$ beta-tensor norm in the actual diagnostic.
For reproducibility the evolved object is the symmetric real-scalar
Weyl tensor $Y_a$, with $S_Y=\sum_bY_b^\dagger Y_b$:
\begin{align}
 \B_{Y_a}&=\tfrac12(S_Y^TY_a+Y_aS_Y)
 +2\sum_bY_bY_a^\dagger Y_b\nonumber\\
 &\quad+\sum_bY_b\operatorname{Re}\tr(Y_b^\dagger Y_a)
 -3\sum_i g_i^2(C_i^TY_a+Y_aC_i). \label{eq:genericbeta}
\end{align}
Projecting \eqref{eq:genericbeta} onto all six actual intertwiners
defines the closed family flow without presuming a holomorphic-pair
cancellation. The scalar and Weyl ordering is unchanged from the
248-real-parent and 48-Weyl common basis.

\section{A genuine momentum-dependent scalar matching subset}
For $V\supset\eta_i(M^2_{ij}+T_{aij}q_a)\eta_j/2$,
the real-scalar functional determinant yields
\begin{equation}
 \Gamma_1=\tfrac12\tr\log(K_0+Tq+\cdots),\qquad
 \Gamma_{1,qq}=-\tfrac14\tr(GTqGTq)+\cdots.
\end{equation}
The UV-finite Euclidean two-point difference is
\begin{multline}
 \Pi_{ab}(p_E^2)-\Pi_{ab}(0)=\frac1{32\pi^2}
 \sum_{ij}^{\text{at least one hard}} T_{aij}T_{bji}\\
 \times\int_0^1\log\left(1+
 \frac{t(1-t)p_E^2}{(1-t)m_i^2+tm_j^2}\right)\dd t.
\end{multline}
Both hard-hard and hard-soft pairs are required. Purely soft pairs
belong to the EFT and are subtracted. Differentiation gives
\begin{align}
 K_{ab}^{S}&=\frac1{32\pi^2}\sum_{ij}^{\rm hard}
 T_{aij}T_{bji}I(m_i^2,m_j^2),\\
 I(x,y)&=\frac{x^2-y^2-2xy\log(x/y)}{2(x-y)^3},&
 I(x,x)&=\frac1{6x},&I(x,0)&=\frac1{2x}.
\end{align}
This is positive semidefinite as a weighted Gram matrix. The single
real-heavy test $ghS^2/2$ gives $K_h=g^2/(192\pi^2M^2)$.
Quartic tadpoles are momentum independent. One-loop mass counterterms
inside this one-loop graph would be selected two-loop contributions.
The massless Goldstone propagators inherit background-field Landau
$\xi=0$; this scalar subset is not separately gauge independent.
With exactly zero soft masses, mixed-bubble soft-region momentum
expansions are scaleless in dimensional regularization. Finite soft
masses or a staged subtraction require a fresh EFT calculation.

All 16 real components of the four doublets are assembled using
the actual scalar Hessian jets and electroweak rotations. The result
obeys all four electroweak Ward identities. Its neutral complex matrix is
\begin{equation}
 K_4^S\simeq\begin{pmatrix}
 .00206838&-.00033316&-.00000208&-.00009074\\
 -.00033316&.00213591&-.00001766&.00000379\\
 -.00000208&-.00001766&.07100660&-.00045735\\
 -.00009074&.00000379&-.00045735&.08599460
 \end{pmatrix},
\end{equation}
with imaginary norm $3.5\times10^{-18}$, checked rather than imposed.
On the bosonic-improved light vector,
\begin{equation}
 \delta Z_h=c_B^\dagger K_4^Sc_B=0.00243676715822.
\end{equation}
The covariant quadratic operator is $(D_\mu H)^\dagger K_4^S(D^\mu H)$.
Gauge vertices required by this operator are not a substitute for
independent gauge-loop, ghost and Nielsen calculations.

For $Z=1+K_4^S$, the canonical operation is
\begin{equation}
 D_c=Z^{-1/2}DZ^{-1/2},\quad
 c_c=\frac{Z^{1/2}c}{\sqrt{c^\dagger Zc}},\quad
 Y_c(c_c)=\frac{Y(c)}{\sqrt{c^\dagger Zc}}.
\end{equation}
The down/lepton expressions use conjugate scalar factors. The zero
eigenvector survives by congruence; its copy-coordinate changes are
not an additional physical light-state rotation. The scalar-leg-only
factor is $0.998783838597$. Exact square roots verify the algebra,
not a two-loop resummation claim. No 55-upper contribution is added
again to this one-step 290-hard calculation.

\section{Mass-ordered seesaw as a changing effective theory}
\subsection{Conventions and running}
Let $H^0=v/\sqrt2$, $V=-m^2H^\dagger H+\lambda(H^\dagger H)^2$ and
\begin{equation}
 M_\nu=\frac{v^2}{2}(C_5-Q),\qquad Q=Y_\nu M_R^{-1}Y_\nu^T.
 \label{eq:Cconvention}
\end{equation}
Store $\widehat M=M_R/\omega$, $\widehat C_5=\omega C_5$.
The actual local type-II input obeys
$\widehat C_5^{II}=2C_{II,\rm raw}f_{\rm raw}$; this factor is tested
against the previous full local I+II matrix. The diagnostic boundary
$g=(.53,.55,.60)$, $\lambda=.12$, $\mu_{\rm start}/\omega=.1$ is
explicitly not a solved P54 matching condition.

Define $H_x=Y_xY_x^\dagger$, $S_N=Y_\nu^\dagger Y_\nu$,
$T=\tr(3H_u+3H_d+H_e+H_\nu)$,
$T_4=\tr(3H_u^2+3H_d^2+H_e^2+H_\nu^2)$ and
$\B=16\pi^2\dd/\dd\log\mu$. In GUT hypercharge normalization,
the $C_{HN}=C_{BN}=0$ subsystem is
\begin{align}
 \B_{Y_u}&=[\tfrac32(H_u-H_d)+T-\tfrac{17}{20}g_1^2
 -\tfrac94g_2^2-8g_3^2]Y_u,\\
 \B_{Y_d}&=[\tfrac32(H_d-H_u)+T-\tfrac14g_1^2
 -\tfrac94g_2^2-8g_3^2]Y_d,\\
 \B_{Y_e}&=[\tfrac32(H_e-H_\nu)+T-\tfrac94g_1^2
 -\tfrac94g_2^2]Y_e,\\
 \B_{Y_\nu}&=[\tfrac32(H_\nu-H_e)+T-\tfrac9{20}g_1^2
 -\tfrac94g_2^2]Y_\nu+3C_5Y_\nu^*M_R,\\
 \B_{M_R}&=S_N^TM_R+M_RS_N,\qquad
 \B_{g_i}=b_ig_i^3,\quad b=(41/10,-19/6,-7),\\
 \B_{C_5}&=P_CC_5+C_5P_C^T+\alpha_CC_5,\\
 P_C&=-\tfrac32H_e+\tfrac72H_\nu,\qquad
 \alpha_C=4\lambda-3g_2^2+2T,\\
 \B_\lambda&=24\lambda^2-3\lambda(3g_2^2+\tfrac35g_1^2)
 +\tfrac34g_2^4+\tfrac38(g_2^2+\tfrac35g_1^2)^2\nonumber\\
 &\quad+4\lambda T-2T_4
 -4\operatorname{Re}\tr(C_5Y_\nu^*M_RY_\nu^\dagger).
\end{align}
The $7/2$, $C_5$ contribution to $\B_{Y_\nu}$ and the last quartic
term include the 2024 correction \cite{Zhang2024}. Our $C_5$ is the
negative of that reference's coefficient, accounting for the two signs.
The former $1/2$ equations are retained only as a negative regression.
This distinction changes tensor trajectories even when its numerical
effect in one chosen benchmark is small.

By the inverse-matrix product rule, with
$P=-3H_e/2+H_\nu/2$,
\begin{align}
 \B_Q&=PQ+QP^T+\alpha_QQ+3(H_\nu C_5+C_5H_\nu^T),\\
 \alpha_Q&=2T-\tfrac9{10}g_1^2-\tfrac92g_2^2,\\
 \B_{C_5-Q}&=P(C_5-Q)+(C_5-Q)P^T+\alpha_C(C_5-Q)
 +(4\lambda+\tfrac9{10}g_1^2+\tfrac32g_2^2)Q. \label{eq:sumflow}
\end{align}
The new terms cancel instantaneously in the difference, but change
$Y_\nu$ and $\lambda$ and thus its later trajectory. Evolving one
initial Weinberg sum above all thresholds would omit the source in
\eqref{eq:sumflow}. The separate II tag is a source decomposition
along one common trajectory, not a new independently fitted matrix.

\subsection{Exact block-tree matching and live events}
At a sterile split let
\begin{equation}
 M'=U^TM_RU=\begin{pmatrix}A&B\\B^T&D\end{pmatrix},\qquad
 Y'=Y_\nu U=(Y_r,Y_h).
\end{equation}
Eliminating $N_h=-D^{-1}(B^TN_r+Y_h^T\ell H)$ gives
\begin{equation}
 \boxed{\ C_5^-=C_5^+-Y_hD^{-1}Y_h^T,\quad
 Y_r^-=Y_r-Y_hD^{-1}B^T,\quad M_r^-=A-BD^{-1}B^T.\ }
\end{equation}
The block inverse identity proves continuity of $C_5-Q$ and
associativity of sequential Schur complements. This is exact matrix
algebra at leading seesaw order, not exact finite-$v$ diagonalization.
For a genuinely mixed non-Takagi block, substituting the heavy solution
also induces a retained kinetic metric; the displayed algebra is not
by itself a canonical EFT matching prescription. Physical threshold
events here use a Takagi basis, where $B=0$ and this issue is absent.

The next event solves
\begin{equation}
 t-\log\sigma_{\max}[\widehat M_R(t)]=0,\qquad t=\log(\mu/\omega),
\end{equation}
with a freshly evaluated Takagi spectrum at each event call.
An entire degenerate cluster is removed at once. Neither frozen
eigenvalues nor a fixed initial ordering defines the event. Tests
include complex family covariance, an actual running level crossing,
non-diagonal Schur matching, a degenerate three-state block and a
minimal two-sterile massless-light-neutrino regression.

The two initial local-light examples have moving event scales
$(.04648960,.02861195,.01786944)$ and
$(.06080995,.03932527,.02503881)$ in $\omega$ units, respectively,
in the displayed zero-$C_{HN}$ diagnostic subsystem. Running continues
after each event and below the last one. These are not certified
physical seesaw scales.

\subsection{The operator-completeness question}
The dimension-five sterile EFT also allows
\begin{equation}
 \mathcal O_{HN}=\overline{N^c}N H^\dagger H,\qquad
 \mathcal O_{BN}=\overline{N^c}i\sigma^{\mu\nu}N B_{\mu\nu}.
\end{equation}
The zero-$C_{HN},C_{BN}$ subspace is invariant under the one-loop
equations, but this does not prove the P54 scalar matching lies in it.
For a heavy real neutral scalar vector $s$, write
\begin{equation}
 V=\tfrac12s^TKs+s^TJH^\dagger H+\cdots,\qquad
 \mathcal L_Y\supset-\tfrac12N^TG_As_AN+\mathrm{h.c.}
\end{equation}
Its tree equation $s=-K^{-1}JH^\dagger H+\cdots$ gives, in the
operator convention with no $1/2$ in $C_{HN}\mathcal O_{HN}$,
\begin{equation}
 C_{HN}=\tfrac12\sum_A G_A(K^{-1}J)_A,\qquad
 \frac{\partial M_R(H)}{\partial(H^\dagger H)}=-2C_{HN}.
 \label{eq:CHNtree}
\end{equation}
Thus it must be computed from the actual neutral scalar source and
Majorana intertwiner. Gauge and PQ zero modes must not be inverted.
This is an operator-basis issue; declaring the coefficient small or
zero without its matching calculation would not repair it.

That calculation is now performed on the exact tree light ray. The
five-real-dimensional SM-singlet scalar plane contains three positive
radial modes and two gauge/PQ zero modes. Inverting only the positive
block in \eqref{eq:CHNtree} gives
\begin{equation}
 \boxed{\quad \omega C_{HN}=-0.0848692011507\,i f_{\rm raw}\ne0.\quad}
\end{equation}
The actual common-phase Majorana intertwiner is used, and its vacuum
contraction reproduces $M_R=i4\sqrt2\sigma f_{\rm raw}$.
An independent check restricts the original potential to the three
radial modes, solves its stationary equations at three nonzero Higgs
fields and extracts $-\Delta M_R/h^2$. It converges to the displayed
coefficient. This is not an assumption inferred from allowed operators.
On the bosonic-improved ray the frozen-tree projection is
$-0.252600617\,i f_{\rm raw}/\omega$; the two local-light projections
are $-0.255045272\,i f_{\rm raw}/\omega$ and
$-0.256536567\,i f_{\rm raw}/\omega$. Those are explicitly mixed-order
projections, not loop-complete Wilson coefficients. The exact-tree
result already disproves a zero-$C_{HN}$ P54 boundary.

\subsection{Required nonzero-$C_{HN}$ evolution and threshold transport}
Keep $C=C_{HN}$, $C_{BN}=0$, and the already declared
$V=-m^2H^\dagger H+\lambda(H^\dagger H)^2$. The additional terms are
\begin{align}
 \Delta\B_{Y_\nu}&=-4Y_\nu M_R^\dagger C,&
 \Delta\B_{M_R}&=+8m^2C,\\
 \Delta\B_\lambda&=16\operatorname{Re}\tr(C^\dagger M_RS_N),\label{eq:CHNlambda}\\
 \B_C&=a_H C+4(CS_N+S_N^TC),&
 a_H&=2T-\tfrac9{10}g_1^2-\tfrac92g_2^2+12\lambda,\\
 \B_{m^2}&=a_Hm^2+4\tr(M_R^\dagger M_RS_N)
 -8\operatorname{Re}\tr(M_R^\dagger M_RM_R^\dagger C).
 \label{eq:CHNflow}
\end{align}
The sign of the last mass equations follows because the cited reference
uses $V=+m_Z^2H^\dagger H+\cdots$, so $m_Z^2=-m^2$ here.
The implementation stores $q_H=m_Z^2/\omega^2=-m^2/\omega^2$
directly, together with $M_R/\omega$ and $\omega C$; this explicit
conversion avoids inferring a sign from the label ``Higgs mass.''
$C_5$ has no direct $C$ mixing at this order. The zero dipole
subsystem is one-loop invariant, and this nonderivative scalar tree
exchange does not generate it. A finite loop dipole coefficient has
not thereby been proven zero.

An independent normalization check uses one real family and
$\mathcal M(h)=\bigl(\begin{smallmatrix}0&yh/\sqrt2\\
yh/\sqrt2&M-Ch^2\end{smallmatrix}\bigr)$.
The coefficient of $h^4$ in $\tr[(\mathcal M^\dagger\mathcal M)^2]$
is $y^4/2-4y^2MC+O(C^2)$. The same Coleman--Weinberg divergence
that gives $\B_\lambda=-2y^4$ therefore gives $+16y^2MC$,
independently checking the coefficient and sign in \eqref{eq:CHNlambda}.

For a non-diagonal sterile split, expand
$M_R(\rho)=M_R-2C\rho$ with $\rho=H^\dagger H$ before taking its
Schur complement. Differentiating once yields
\begin{multline}
 C^-=C_{rr}-C_{rh}D^{-1}B^T-BD^{-1}C_{hr}\\
 +BD^{-1}C_{hh}D^{-1}B^T. \label{eq:CHNSchur}
\end{multline}
This reduces to the retained block in an exact Takagi basis. The
$\rho$-dependent Yukawa shift and Weinberg denominator produce operators
of dimension six and seven, respectively, outside this linear
dimension-five truncation. They must not be mistaken for additional
dimension-five matching terms.

Equation \eqref{eq:sumflow} is expressly the zero-$C$ identity. With
nonzero $C$, its additional source is minus
\begin{equation}
 \Delta\B_Q=(\Delta\B_{Y_\nu})M_R^{-1}Y_\nu^T
 +Y_\nu M_R^{-1}(\Delta\B_{Y_\nu})^T
 -Y_\nu M_R^{-1}(\Delta\B_{M_R})M_R^{-1}Y_\nu^T.
\end{equation}
Thus a full model run cannot continue to assert the zero-$C$ sum-flow
identity after the nonzero boundary is inserted. A diagnostic choice
$m^2=0$ at one starting scale is not radiatively invariant and is not
a physical electroweak matching condition.
In the two nonzero-$C$ diagnostic runs the terminal values are
$q_H=2.22133\times10^{-6}$ and $2.98359\times10^{-6}$, respectively.
They exceed the squared terminal scale in $\omega$ units: a Higgs
kept light all the way to that endpoint is not self-consistent without
quadratic matching and retuning. The mass-independent $\MS$ ODE can
be continued as a mathematical diagnostic, but its fixed-$v$ terminal
mass proxies are not a physically valid light-Higgs EFT endpoint.

\section{Reproducibility and remaining physical work}
The code and JSON/Markdown records under \texttt{route\_f/code} and
\texttt{route\_f/output} are named
\begin{center}\small
\begin{tabular}{ll}
\toprule
Stem after \texttt{verify\_} & Computed scope\\
\midrule
\texttt{p54\_lower\_covariant\_eft} & scalar valley, vector response, local vertices\\
\texttt{p54\_upper\_yukawa\_thresholds} & actual pure-Yukawa finite state sums\\
\texttt{p54\_scalar\_kinetic\_matching} & full scalar-cubic four-doublet metric\\
\texttt{p54\_scalar\_chn} & actual neutral-scalar tree $NNH^\dagger H$ coefficient\\
\texttt{p54\_sequential\_seesaw} & moving thresholds and dimension-five flow\\
\bottomrule
\end{tabular}
\end{center}
Each record declares its boundary data, source hashes, conventions,
scope restrictions and executable positive and negative tests.
Previously computed expensive Hessians are reused.

The next physical matching calculation must include the omitted
heavy-vector fermion kinetic/vertex graphs, the vector/scalar/Goldstone/
ghost scalar kinetic package, upper scalar matching and tadpole
responses, and the lower covariant hard-minus-EFT subtraction.
Sequential finite one-loop type-I/type-II matching is separate from
the tree Schur rule: an all-sterile formula cannot silently be reused
when lighter sterile states and an existing type-II operator remain.
Only after these matched inputs are assembled is a constrained global
fit meaningful. Momentum-dependent pole self-energies and Nielsen
consistency remain separate from a zero-momentum curvature condition.
No soliton regularity result, determinant or portal is promoted here.
The sterile running engine is a specified $C_{BN}=0$ dimension-five
subsystem with the PQ axion held as a background spectator. An invariant
dipole-free RGE subspace is not a proof of zero finite dipole matching,
and freezing the axion background is not a calculation of quantum axion
loops. Its Wilson couplings and quantum effects must be accounted for
before calling the result the full physical P54 lower evolution.

\begin{thebibliography}{9}
\bibitem{Zhang2024} D. Zhang, \emph{Threshold Effects on the Massless
Neutrino in the Canonical Seesaw Mechanism} (2024),
\url{https://arxiv.org/abs/2405.18017}.
\bibitem{Antusch} S. Antusch et al., \emph{Running Neutrino Mass
Parameters in See-Saw Scenarios} (2005),
\url{https://arxiv.org/abs/hep-ph/0501272}.
\bibitem{Patel} K. M. Patel and V. Shukla, finite Yukawa matching
conventions and kernels, \url{https://arxiv.org/abs/2310.16563}.
\bibitem{Luo} M. Luo, H. Wang and Y. Xiao, general renormalization-group
equations, \url{https://arxiv.org/abs/hep-ph/0211440}.
\bibitem{Matching} M. Gabelmann, M. Muehlleitner and F. Staub,
\emph{Automatised matching between two scalar sectors at the one-loop level},
\url{https://doi.org/10.1140/epjc/s10052-019-6570-5}.
\bibitem{TypeI} D. Zhang and S. Zhou, complete one-loop type-I matching,
\url{https://arxiv.org/abs/2107.12133}.
\end{thebibliography}
\end{document}
